% Supplementary material — single-column, shares zhconf typography.
% Build: make -C paper zh_conf_supp
%   or:  cd paper/zh_conf && xelatex supp.tex && bibtex supp && xelatex supp.tex && xelatex supp.tex
%
% Design notes:
%   - Separate PDF from main.tex (standard for conference supplementary material).
%   - Single column (onecolumn): the null-space proofs in App. A and the wide
%     benchmark tables in App. D do not fit comfortably in the 2-column main text.
%   - Reuses zhconf.sty for identical fonts/citation style; geometry widened for
%     single-column reading width.

\documentclass[onecolumn,11pt,a4paper]{article}
\usepackage{zhconf}

% Single-column reading width (override the two-column geometry from zhconf).
\geometry{left=2.6cm,right=2.6cm,top=2.4cm,bottom=2.4cm}

\usepackage{longtable}

% 长代码路径较多，放宽断行避免 overfull（允许稍宽的词间距而非溢出页边）。
\sloppy

% Appendix-style section numbering: A, A.1, A.1.1 ...
\renewcommand{\thesection}{\Alph{section}}
\renewcommand{\thesubsection}{\thesection.\arabic{subsection}}
\renewcommand{\thesubsubsection}{\thesubsection.\arabic{subsubsection}}

% Show subsubsections in the table of contents.
\setcounter{secnumdepth}{3}
\setcounter{tocdepth}{2}

% Proposition lead-in helper: \prop{P1}{精确零空间不变性}
\newcommand{\prop}[2]{\medskip\noindent\textbf{命题~#1（#2）.}\ }
\newcommand{\lemmastate}[2]{\medskip\noindent\textbf{引理~#1（#2）.}\ }
\newcommand{\coro}[2]{\medskip\noindent\textbf{推论~#1（#2）.}\ }

\title{%
  \zhconftitle{%
    补充材料\\
    \large 细节更多，却在漂移：无真值热成像多帧超分辨中的零空间幻觉
  }%
}
\author{\zhconfauthor{匿名作者（投稿前替换）}}
\date{}

\begin{document}

\maketitle

\noindent\small
本补充材料是同名主文的附录，单栏排版，编号自成体系（A--E）。文中所有「主文~\S」指向正文章节，
「附录~A.2」「见~D.4.3」等指向本文件内部。本文件与主文共享同一参考文献库与图表资产，
不重复正文已给出的结论，只补充推导、协议、完整表格与复现细节。
\normalsize

\medskip
\hrule
\medskip

\tableofcontents

\bigskip

\section{理论与推导}

\subsection{系统 MTF 与 SR 可行性分析}

本节从系统的调制传递函数（\eng{MTF}）和信噪比（\eng{SNR}）出发，建立 \(2\times\) 与 \(4\times\) 超分辨率的理论可行性边界，
为主文~\S3.1 和~\S1 的 claim 边界提供定量支撑。

\subsubsection{系统 MTF 模型}

本系统的光学 \eng{PSF} 建模为各向同性 Gaussian，标准差 \(\sigma\) 以 \eng{LR} 像素为单位（1 \eng{LR} px \(=\) 10\,\(\mu\)m）。对应的 \eng{MTF} 为
\[
  \mathrm{MTF}_{\mathrm{PSF}}(f) = \exp\!\bigl(-2\pi^2 \sigma^2 f^2\bigr),
\]
其中 \(f\) 为空间频率（cyc/\eng{LR} px）。探测器孔径的 \eng{MTF} 为
\[
  \mathrm{MTF}_{\mathrm{det}}(f) = \bigl|\mathrm{sinc}(W f)\bigr|, \quad W = 10\;\mu\mathrm{m}.
\]
系统总 \eng{MTF} 为二者之积。\eng{EP15 M3} 的 \eng{FRC} 形状拟合使用了 \(\mathrm{MTF}^2\) 形式，即
\[
  \mathrm{MTF}(f)^2 = \bigl[\exp\!\bigl(-2\pi^2\sigma_{\mu\mathrm{m}}^2 f^2\bigr) \cdot \bigl|\mathrm{sinc}(10f)\bigr|\bigr]^2, \quad \sigma_{\mu\mathrm{m}} = \sigma_{\mathrm{LR}} \times 10\;\mu\mathrm{m}.
\]
需要注意的是，主文~\S3.1 引用的 \eng{MTF} 数值表仅含 Gaussian \eng{PSF} 项，未乘 detector sinc；\eng{EP15 M3} 拟合使用的则是含 sinc 的系统 \(\mathrm{MTF}^2\)。下文分别标明各自口径。

当前系统的已校准空间分辨率为 20\,\(\mu\)m，探测器采样间距（pitch）为 10\,\(\mu\)m/px，对应 \eng{FOV} 6.4\(\times\)4.8\,mm。
在 \(1\times/2\times/4\times\) 输出网格上，Nyquist 频率分别为 0.5、1.0、2.0\,cyc/detector px。下表给出 Gaussian-only 口径下各网格 Nyquist 处的 \eng{MTF} 值。

\begin{table}[ht]
  \centering
  \caption{各输出网格 Nyquist 频率处的 Gaussian-only \eng{MTF} 值}
  \label{tab:mtf-nyquist}
  \small
  \begin{tabular}{lcccc}
    \toprule
    网格 & \(f\) (cyc/px) & \(\sigma=0.20\) & \(\sigma=0.35\) & \(\sigma=0.50\) \\
    \midrule
    \(1\times\) Nyquist & 0.5 & 0.821 & 0.546 & 0.291 \\
    \(2\times\) Nyquist & 1.0 & 0.454 & 0.089 & 0.007 \\
    \(4\times\) Nyquist & 2.0 & 0.042 & \(6.3\times10^{-5}\) & \({\sim}3\times10^{-9}\) \\
    \bottomrule
  \end{tabular}
\end{table}

\subsubsection{有效 SNR 判据}

为判断某一空间频率处的信息是否可恢复，定义有效 \eng{SNR} 为
\[
  \mathrm{SNR}_{\mathrm{eff}}(f) = \frac{\Delta T \cdot \mathrm{MTF}(f,\sigma)}{\sigma_n},
\]
其中 \(\Delta T\) 为热对比度，\(\sigma_n = 0.0724\;^\circ\mathrm{C}\) 为探测器噪声底（由平滑区域相邻坐标 \eng{MAE} 测定）。
实测对比度覆盖从噪声底到外轮廓中位 2.49\,\(^\circ\)C 的宽范围：名义边缘 0.70\,\(^\circ\)C（输入 \eng{SNR} \(=\) 9.7）、内轮廓中位 1.94\,\(^\circ\)C（\eng{SNR} \(=\) 26.8）、外轮廓中位 2.49\,\(^\circ\)C（\eng{SNR} \(=\) 34.4）。

将上述对比度代入各网格 Nyquist 处 \eng{MTF}，得到有效 \eng{SNR} 表。

\begin{table}[ht]
  \centering
  \caption{各对比度在 \(2\times\)/\(4\times\) Nyquist 处的有效 \eng{SNR}}
  \label{tab:snr-eff}
  \small
  \begin{tabular}{lcc}
    \toprule
    对比度 & \(2\times\)（\(\sigma\)=0.20 / 0.35 / 0.50） & \(4\times\)（\(\sigma\)=0.20 / 0.35 / 0.50） \\
    \midrule
    名义 0.7\,\(^\circ\)C    & 4.39 / 0.86 / 0.07   & 0.41 / 0.001 / \(\sim\)0 \\
    内轮廓 1.94\,\(^\circ\)C & 12.16 / 2.39 / 0.19  & 1.14 / 0.002 / \(\sim\)0 \\
    外轮廓 2.49\,\(^\circ\)C & 15.62 / 3.06 / 0.25  & 1.46 / 0.002 / \(\sim\)0 \\
    \bottomrule
  \end{tabular}
\end{table}

风险带定义为：\(\mathrm{SNR}_{\mathrm{eff}} \ge 5\) 为 observable，\(\ge 3\) 为 borderline，\(\ge 1\) 为 weak，其余为 noise-dominated。

\subsubsection{可行性判定}

\textbf{\(2\times\) 在条件下可行。} 在乐观 \(\sigma\) 端（\(\sigma \le 0.35\)），内/外轮廓的 \(\mathrm{SNR}_{\mathrm{eff}}\) 处于 observable 到 borderline 区间，
存在可用的频率余量。但当 \(\sigma = 0.5\) 时，\(2\times\) Nyquist 处 \eng{MTF} 仅剩 0.007，高频信号强烈衰减至噪声以下。
因此本 \eng{POC} 将目标定位在 contour-level 增强，而非计量级温度恢复。

\textbf{\(4\times\) 出界。} 除最乐观的 \(\sigma = 0.2\) 在外轮廓对比度下勉强达到 weak（1.46），其余条件下 \(\mathrm{SNR}_{\mathrm{eff}}\) 全线 noise-dominated。
\(4\times\) 网格仅作 contour oversampling 和可视化用途。这一理论预测与后续 \(4\times\) 网络实验的负结果一致（见 D.4.2）——理论界先于实验设定了风险预期。

需要强调的是，\(\mathrm{SNR}_{\mathrm{eff}}\) 过门是必要条件而非充分证明：该判据忽略了对齐误差、热漂移与模型失配等因素。

作为辅助证据，Cramér-Rao 下界（\eng{CRB}）给出了边缘定位精度的理论极限。例如在 \(\Delta T = 0.7\;^\circ\mathrm{C}\)、\(\sigma_{\mathrm{PSF}} = 0.5\;\mathrm{px}\) 条件下，
单帧 \eng{CRB} 为 0.1415\,px，16 帧平均后降至 0.0344\,px。\eng{EP04} 实测 A-class split-half 约 0.027\,px 与 16 帧 \eng{CRB} 同量级，二者作为一致性检验互相支撑。

\subsection{观测算子与零空间}

\subsubsection{前向模型与记号}

主文~\S3.1 定义的离散观测模型为
\[
  y_k = D \cdot B \cdot H \cdot S_{t_k} \cdot x + n_k,
\]
其中 \(x\) 为 \(2\times\) 网格上的 \eng{HR} 图像，\(S_{t_k}\) 以 contour-refined 位移 \(t_k\) 对 \eng{HR} 网格作亚像素平移重采样，
\(H\) 为 Gaussian \eng{PSF} 卷积（\(\sigma \in [0.2, 0.5]\) \eng{LR} px，见 A.5.2），\(B\) 为 10\,\(\mu\)m 探测器孔径 box 积分，
\(D\) 为 \eng{HR}\(\rightarrow\)\eng{LR} 降采（scale\,=\,2），\(n_k\) 为噪声底 0.0724\,\(^\circ\)C。

仓库中有两套该模型的实现：\eng{EP06} 的矩阵式 forward/adjoint 用于经典重建（\eng{MAP-TV}/\eng{TGV}），
\eng{EP07} 训练 loss 链（blur \(\rightarrow\) block-average \(\rightarrow\) 可选带限）用于深度学习。两者共享 \(H\) 与 \(B \cdot D\) 的核心低通-折叠结构。下文的零空间论证对两者同样成立。

固定 \(t_k\) 后，每帧算子
\[
  A_k = D \cdot B \cdot H \cdot S_{t_k}
\]
是线性的；248 帧 clean 主 session 的 stacked burst 算子写为
\[
  A = \begin{bmatrix} A_1 \\ \vdots \\ A_K \end{bmatrix}, \qquad K=248.
\]
除特别说明外，\(\|\cdot\|\) 为 Euclidean norm。令 \(A = U\Sigma V^\top\) 为完整奇异值分解，右奇异向量为 \(\{v_i\}\)，奇异值为
\[
  \sigma_1 \ge \sigma_2 \ge \cdots \ge \sigma_N \ge 0,\qquad A v_i = \sigma_i u_i.
\]
精确零空间与 \(\epsilon\)-零空间定义为
\[
  \ker(A)=\mathrm{span}\{v_i:\sigma_i=0\},\qquad
  \mathcal N_\epsilon=\mathrm{span}\{v_i:\sigma_i\le\epsilon\}.
\]
观测域损失统一写作 \(\mathcal L(x)=\ell(Ax,y)\)。
若使用 highpass/full-band 等观测侧频带锚，则写作 \(\mathcal L_W(x)=\ell(WAx,Wy)\)，其中 \(W\) 是作用在观测域的有界线性算子。
下列命题只证明固定线性观测算子的性质：它们不证明网络或 \eng{SGD} 一定漂移；它们证明的是，一旦某个漂移分量落在 \(\ker(A)\) 或 \(\mathcal N_\epsilon\) 内，
有界观测侧锚定就只能提供零或趋零的信息与恢复力。

\subsubsection{零空间刻画}

令 \(A = [A_1;\ldots;A_K]\)。stacked 算子的零空间由逐帧零空间交集给出：
\[
  \ker(A)=\bigcap_{k=1}^K \ker(A_k),
\]
且对任意扰动 \(\delta x\)，
\[
  \|A\delta x\|^2 = \sum_{k=1}^K \|A_k\delta x\|^2.
\]
因此，一个单位方向对整个 burst 近似不可观测，当且仅当它在所有帧上的 root-sum-square 响应都很小。
多帧相位多样性只能增加约束：它把单帧条件替换为 \(K\) 个条件的交集，从而可缩小 exact null space，并改善覆盖良好的 alias block 的奇异值；
但它不能恢复被共同的 pre-sampling transfer function \(B H\) 消掉或压到噪声以下的频率。

从机制上看，零空间源于两个来源的叠加。

\textbf{第一，带限衰减。} \(H\) 和 \(B\) 均为低通算子——如 A.1.1 所示，Gaussian \eng{MTF} 在 \(2\times\) Nyquist 处最低仅剩 0.007（\(\sigma = 0.5\)），
叠加 box sinc 后更低。\eng{HR} 网格上超过截止频率的高频分量经 \(A\) 映射后衰减至噪声以下，在数值意义上落入 \(\epsilon\)-零空间。

\textbf{第二，混叠折叠。} 降采算子 \(D\) 将 \eng{HR} 频谱折叠进 \eng{LR} Nyquist 带宽，存在整族高频扰动 \(\delta x\) 使得折叠后各别名分量相消，
严格满足 \(A \cdot \delta x = 0\)。

在周期、shift-invariant 的近似下，这一点可写成 Fourier alias block。对 \(2\times\) 降采，一个 \eng{LR} base frequency \(\omega\) 接收四个 \eng{HR} alias：
\[
  \mathcal A(\omega)=\{\omega+\pi q:q\in\{0,1\}^2\}.
\]
若 \(m(\xi)\) 是 \(B H\) 在 \eng{HR} 频率 \(\xi\) 处的传递系数，则第 \(k\) 帧观测系数为
\[
  \widehat{(A_k\delta x)}(\omega)
  =\sum_{\xi\in\mathcal A(\omega)} m(\xi)e^{-i\xi\cdot t_k}c_\xi.
\]
对应的 stacked alias block 为
\[
  \Phi_\omega(k,\xi)=m(\xi)e^{-i\xi\cdot t_k}.
\]
精确 alias-cancellation 方向就是 \(\ker(\Phi_\omega)\)。更多相位可提高 \(\Phi_\omega\) 的 rank，去掉部分 downsampling-only alias nulls；
但若某个 \(m(\xi)=0\)，该频率是所有帧共有的 transfer zero，任意相位都不可见。若 \(m(\xi)\) 非零但很小，则该方向进入 \(\epsilon\)-零空间，统计上仍不可靠。

直观地说，「细于截止周期的结构改动」与「在所有已采相位下相消的别名组合」对 \(1\times\) 观测不可见或近似不可见。
多帧相位多样性可缩小但不消灭该盲区（25 个相位 bin 的覆盖分析见 A.4.1 和 B.2.3）。

\subsubsection{P1：精确零空间不变性}

\prop{P1}{精确零空间不变性}设 \(\delta x_{\mathrm{null}}\) 满足 \(A \cdot \delta x_{\mathrm{null}} = 0\)，
\(\mathcal{L}(\hat{x}) = \ell(A\hat{x}, y)\) 为任意仅经 \(A\) 接触预测的观测域损失。则对任意 \(\hat{x}\)：
\[
  \mathcal{L}(\hat{x} + \delta x_{\mathrm{null}}) = \ell(A\hat{x} + A\delta x_{\mathrm{null}},\; y) = \ell(A\hat{x},\; y) = \mathcal{L}(\hat{x}).
\]
若 \(\ell\) 对第一自变量可微，则链式法则给出 \(\nabla\mathcal L(\hat{x})=A^\top\nabla_1\ell(A\hat{x},y)\)，因此
\[
  \langle \nabla\mathcal L(\hat{x}),\delta x_{\mathrm{null}}\rangle
  =\langle \nabla_1\ell(A\hat{x},y),A\delta x_{\mathrm{null}}\rangle=0.
\]
此命题的推论是：对 exact null component，任意有限权重的 forward-consistency loss、任意有界 highpass/full-band 频带权重、
以及任意只通过 \(A\) 接触预测的残差范数，都没有可用信息。这正是主文~\S5.2 中 \eng{Stats+HP-FC} / \eng{Stats+Full-FC} 与无锚 \eng{Stats}（run ids: V9B/V9D/v8.1a）漂移曲线重合的机制解释（数字见 D.4.3）。

\subsubsection{P2：\texorpdfstring{\(\epsilon\)}{epsilon}-零空间灵敏度界}

\prop{P2}{\(\epsilon\)-零空间灵敏度}对任意扰动 \(\delta x\)，
\[
  \|A\delta x\|^2=\sum_i \sigma_i^2\langle v_i,\delta x\rangle^2.
\]
若 \(\ell\) 在 \(A\hat{x}\) 与 \(A(\hat{x}+\delta x)\) 之间的线段上对第一自变量是 \(L_\ell\)-Lipschitz，则
\[
  |\mathcal L(\hat{x}+\delta x)-\mathcal L(\hat{x})|
  \le L_\ell\|A\delta x\|.
\]
特别地，若 \(\|\delta x\|=1\) 且 \(\delta x\in\mathcal N_\epsilon\)，则
\[
  \|A\delta x\|\le\epsilon,\qquad
  |\mathcal L(\hat{x}+\delta x)-\mathcal L(\hat{x})|\le L_\ell\epsilon.
\]
对带权观测域损失 \(\mathcal L_W(\hat{x})=\ell(WA\hat{x},Wy)\)，若 \(W\) 有界，则
\[
  |\mathcal L_W(\hat{x}+\delta x)-\mathcal L_W(\hat{x})|
  \le L_\ell\|W\|\,\|A\delta x\|.
\]
因此，\(\|W\|\le1\) 的 normalized highpass/full-band 权重不会比未加权观测损失更敏感；未归一化但有界的 \(W\) 也只能把界放大有限倍。
若 \(A\delta x\) 恰落在 \(W\) 抑制的频带内，灵敏度还会进一步下降，甚至为零。

证明只需把 \(\delta x\) 展开到右奇异向量基中：
\[
  \delta x=\sum_i\langle v_i,\delta x\rangle v_i,\qquad
  A\delta x=\sum_i\sigma_i\langle v_i,\delta x\rangle u_i,
\]
再应用 Lipschitz 不等式与 \(\|WA\delta x\|\le\|W\|\|A\delta x\|\)。
注意，对平方损失而言 Lipschitz 常数是局部的；下一命题直接计算平方损失，不依赖全局 Lipschitz 假设。

\subsubsection{P3：平方损失恢复曲率与有限权重推论}

\prop{P3}{平方观测损失的恢复曲率}对平方数据项
\[
  \mathcal L(\hat{x})=\frac12\|A\hat{x}-y\|^2,
\]
有 \(\nabla\mathcal L(\hat{x})=A^\top(A\hat{x}-y)\)。沿右奇异方向 \(v_i\)，
\[
  \langle\nabla\mathcal L(\hat{x}),v_i\rangle
  =\sigma_i\langle u_i,A\hat{x}-y\rangle
  =\sigma_i^2\langle v_i,\hat{x}\rangle-\sigma_i\langle u_i,y\rangle,
\]
且在直线 \(\hat{x}+\alpha v_i\) 上，
\[
  \frac{d^2}{d\alpha^2}\mathcal L(\hat{x}+\alpha v_i)=\sigma_i^2.
\]
因此，有限权重 \(\lambda\mathcal L\) 在 \(v_i\) 方向上的恢复曲率为 \(\lambda\sigma_i^2\)，随 \(\sigma_i\to0\) 消失。对有界频带加权平方损失
\(\mathcal L_W(\hat{x})=\frac12\|W(A\hat{x}-y)\|^2\)，该方向的曲率为
\[
  \frac{d^2}{d\alpha^2}\mathcal L_W(\hat{x}+\alpha v_i)=\|WAv_i\|^2\le\|W\|^2\sigma_i^2.
\]
所以，有界观测频带权重不能在 \(A\) 的 near-null 方向制造非消失曲率。

\coro{P3}{有限观测权重无法平衡 \(O(1)\) 近零空间力}固定 \(\sigma_i>0\) 的奇异方向 \(v_i\)。
若除有限权重观测项 \(\lambda\mathcal L\) 外，还有一个先验项或网络诱导项 \(G\)，其在 \(v_i\) 方向的局部导数为非零常数 \(g_i\)，
并在所考虑的小段内近似不变，则相对 data-only stationary point 的一维平衡位移必须满足
\[
  |\alpha|=\frac{|g_i|}{\lambda\sigma_i^2}.
\]
若使用带权观测损失，分母变为 \(\lambda\|WAv_i\|^2\)，且 \(\lambda\|WAv_i\|^2\le\lambda\|W\|^2\sigma_i^2\)。
于是，对任意固定有限 \(\lambda\) 与有界 \(W\)，若竞争力 \(g_i=O(1)\)，平衡该力所需的位移会随 \(\sigma_i\to0\) 发散。
反过来说，若人为选择量级为 \(1/\sigma_i^2\) 的无界逆 \eng{MTF} 权重，代数上确实可制造 \(O(1)\) 曲率；
但这不是新的观测信息，而是逆奇异值放大。P4 说明同一放大会同步放大噪声。

\subsubsection{P4：统计不可辨识方差下界}

\prop{P4}{奇异坐标的方差下界}假设
\[
  y=A x^\star+n,\qquad \mathbb E[n]=0,\qquad \mathrm{Cov}(n)=\eta^2 I.
\]
令 \(c_i=\langle v_i,x^\star\rangle\)。对任何线性估计器 \(\hat c_i=a^\top y\)，若它对所有 \(x^\star\) 都无偏估计 \(c_i\)，则
\[
  \mathrm{Var}(\hat c_i)\ge\frac{\eta^2}{\sigma_i^2}.
\]
等号由 \(a=u_i/\sigma_i\) 达到。若 \(\sigma_i=0\)，则不存在对所有 \(x^\star\) 无偏的线性估计器。若噪声为 Gaussian，同一界也是该正交坐标的 Cramér-Rao 下界。

证明要点是，无偏性要求 \(A^\top a=v_i\)。与 \(v_i\) 取内积得
\[
  1=\langle A^\top a,v_i\rangle=\langle a,Av_i\rangle=\sigma_i\langle a,u_i\rangle\le\sigma_i\|a\|,
\]
故 \(\|a\|^2\ge1/\sigma_i^2\)，而 \(\mathrm{Var}(a^\top y)=\eta^2\|a\|^2\)。

P1、P2/P3 与 P4 分别对应三层失效：\(\sigma_i=0\) 时 exact invariance；\(\sigma_i\) 很小时，观测损失变化为 \(O(\sigma_i)\)、平方损失恢复曲率为 \(O(\sigma_i^2)\)；
在带噪声数据中，\(c_i\) 的 Fisher information 也只有 \(\sigma_i^2/\eta^2\)。
因此，near-null 成分必须靠先验、输入证据或选点协议管理，而不能指望另一个有限观测侧锚定项恢复。

\subsubsection{两曲线诊断与直接测量状态}

\textbf{两曲线诊断的充分条件。} 若训练过程中 (i) forward loss 已贴近其噪声地板而无法继续下降，
且 (ii) 真实数据上的 proxy 指标（artifact 与 corr，定义见 A.3）持续单调漂移，
则这强烈指示漂移方向的观测域投影约为零——即漂移主要发生在 \(A\) 的 \(\epsilon\)-零空间内。
该判断是充分条件式的推断，不是恒等式，也不证明网络必然漂移；它只说明一旦实测漂移同时满足这两条曲线特征，有界观测侧锚定缺少可用恢复力。
在 \eng{Stats+HP-FC}（v9b）实验中，forward loss 自 10K 步起平坦于 0.004--0.009，而同期 artifact 从 0.37 单调上爬至 0.65，恰好满足该充分条件。

\textbf{零空间投影的直接测量（待补充 / stretch goal）。} 计划取同一变体相邻 checkpoint 的预测差 \(\delta\hat{x} = \hat{x}_{k+1} - \hat{x}_k\)，
计算 range 分量 \(\delta x_{\mathrm{range}} = A^\dagger A \cdot \delta\hat{x}\)（\(A^\dagger\) 用共轭梯度近似），报告 \(\|\delta x_{\mathrm{range}}\| / \|\delta\hat{x}\|\) 随训练步数的变化。
若该比值在漂移期趋近零，则可将间接诊断升级为直接测量。该项尚未完成，不能作为本文已交付证据。

\subsection{Proxy 反相关的构造性论证}

主文使用 artifact score 与 raw-control correlation 构成的 proxy 对来监测训练漂移。本节论证二者为何存在构造性的反相关关系。

\subsubsection{指标定义}

两个 proxy 共享同一预处理：将 \eng{SR} 输出温度图变换到 highpass 域 \(u = \hat{x} - G_{\sigma_{\mathrm{bg}}}(\hat{x})\)（\(\sigma_{\mathrm{bg}} = 5\) \eng{HR} px），
控制图 \(c\) 为 248 帧原始数据均值的 bicubic 上采样经同一 highpass 变换。

TB-scale 伪影分数 \(\artscore(u)\)（用于训练期 eval 和主文 F3/F4 轨迹图）定义为
\[
  \artscore(u) = \frac{\mathrm{std}(h(u)) + 0.25 \cdot \mathrm{std}(\nabla^2 u)}{\mathrm{std}(u)},\quad h(u)=u-G_{\sigma=1}(u),
\]
即 \(u\) 内部更高频能量的占比。\eng{EP11-harness} 口径使用不同定义（ringing \(+\) 0.25\(\times\)blockiness），数值与 TB-scale 不可比，不得混入同一表格。

Raw-control correlation 定义为 \(\mathrm{corr} = \mathrm{Pearson}(u, c)\)，取全图 finite 像素。

\subsubsection{一阶反相关引理}

两个 proxy 都是同一张 highpass 图 \(u\) 的泛函——artifact 度量 \(u\) 的内部高频能量分布，corr 度量 \(u\) 与固定观测锚 \(c\) 之间的线性一致性。
它们并非独立的证据源，这是反相关论证的前提。但反相关并非无条件成立：只有当扰动比当前残差更高频、且不与固定 raw-control highpass map 对齐时，
artifact 才会一阶上升、corr 才会一阶下降。

为写出条件，令所有向量都已去均值，且 \(c\ne0\)。定义
\[
  L_1=I-G_{\sigma=1},\qquad L_2=\nabla^2,\qquad \beta=0.25,
\]
并令 \(N(u)=\|L_1u\|+\beta\|L_2u\|\)，\(S(u)=\|u\|\)。
于是 artifact proxy 等价于 \(N(u)/S(u)\)（std 的常数因子在比值中抵消），raw-control correlation 为
\[
  \mathrm{corr}(u)=\frac{\langle u,c\rangle}{\|u\|\|c\|}.
\]

\lemmastate{L5}{一阶反相关}令 \(u(\alpha)=u_0+\alpha e\)，并假设 \(\|u_0\|>0\)、\(\|L_1u_0\|>0\)、\(\|L_2u_0\|>0\)，且 \(\rho_0:=\mathrm{corr}(u_0)>0\)。
定义 broadband alignment \(s_0(e)=\dfrac{\langle u_0,e\rangle}{\|u_0\|^2}\)，以及 high-frequency directional growth
\[
  h_0(e)=
  \frac{
  \frac{\langle L_1u_0,L_1e\rangle}{\|L_1u_0\|}
  +\beta\frac{\langle L_2u_0,L_2e\rangle}{\|L_2u_0\|}
  }{
  \|L_1u_0\|+\beta\|L_2u_0\|
  }.
\]
若 \(h_0(e)>s_0(e)\)，则
\[
  \left.\frac{d}{d\alpha}\artscore(u_0+\alpha e)\right|_{\alpha=0}>0.
\]
若进一步满足 \(\dfrac{\langle e,c\rangle}{\|u_0\|\|c\|}<\rho_0\,s_0(e)\)，则
\[
  \left.\frac{d}{d\alpha}\mathrm{corr}(u_0+\alpha e)\right|_{\alpha=0}<0.
\]
也就是说，对边缘增强式的 style perturbation，只要它相对当前 \(u_0\) 增加更高频能量，且与固定 raw-control map \(c\) 的对齐很弱，
artifact 会一阶上升而 corr 会一阶下降。

证明来自范数一阶导数。对任意非零可微路径 \(q(\alpha)\)，\(\left.\frac{d}{d\alpha}\|q(\alpha)\|\right|_{\alpha=0}=\frac{\langle q(0),q'(0)\rangle}{\|q(0)\|}\)。
代入 \(L_1u(\alpha)\)、\(L_2u(\alpha)\) 与 \(u(\alpha)\)，得到
\[
  N'(0)=
  \frac{\langle L_1u_0,L_1e\rangle}{\|L_1u_0\|}
  +\beta\frac{\langle L_2u_0,L_2e\rangle}{\|L_2u_0\|},
  \qquad
  S'(0)=\frac{\langle u_0,e\rangle}{\|u_0\|}.
\]
由 \(\artscore(u)=N(u)/S(u)\) 可知，\(\artscore\) 的一阶导数为正当且仅当 \(N'(0)/N(u_0)>S'(0)/S(u_0)\)，即 \(h_0(e)>s_0(e)\)。对 correlation，
\[
  \left.\frac{d}{d\alpha}\mathrm{corr}(u_0+\alpha e)\right|_{\alpha=0}
  =
  \frac{\langle e,c\rangle}{\|u_0\|\|c\|}
  -\rho_0\frac{\langle u_0,e\rangle}{\|u_0\|^2},
\]
于是上式小于零正是第二个条件。

若 \(e\perp u_0\) 且 \(e\perp c\)，corr 的一阶导数为零而非负。此时若 \(\rho_0>0\)，下降只在二阶出现：
\[
  \mathrm{corr}(u_0+\alpha e)
  =\rho_0\left(1+\alpha^2\frac{\|e\|^2}{\|u_0\|^2}\right)^{-1/2}.
\]
因此，主文中的「构造性反相关」应理解为满足上述方向条件的局部诊断，而不是任意扰动方向上的恒等关系。

由此得到三条推论。第一，proxy 对是漂移温度计与选点准则，不是两个可联合优化的分数。
第二，跨 input-mode 的横向比较无效：hybrid drizzle 输入合法携带更多高频能量，其 artifact 基线天然更高。
第三，应关注训练轨迹的走向而非端点数值。

\subsubsection{假设与适用边界}

\begin{itemize}
  \item \textbf{固定线性算子。} 上述命题假设 \(t_k\) 已固定，因此 \(S_{t_k}\) 与 \(A_k=D B H S_{t_k}\) 是线性算子。若把位移与 \(x\) 联合优化，问题不再是本文证明的固定线性观测算子。
  \item \textbf{有限维离散化。} \eng{SVD} 命题对实现中的有限矩阵精确成立；Fourier/MTF 解释还额外假设周期或内部区域近似 shift-invariant。非周期边界会扰动 Fourier 对角化，但不影响代数 \eng{SVD} 结论。
  \item \textbf{相位多样性。} stacked exact null space 是 \(\bigcap_k\ker(A_k)\)。更多相位可去掉部分单帧混叠零空间，却不能去掉 \(B H\) 的共同 transfer zeros，也不能让极小 \eng{MTF} 方向变得统计可靠。
  \item \textbf{有界观测侧带权。} 「频带锚定不能补救」指的是有界 \(W\) 与固定有限损失权重 \(\lambda\)。无界逆滤波或按 \(1/\sigma_i^2\) 放大的权重可在代数上造曲率，但会按同一逆奇异值放大噪声，不属于有限观测锚定主张。
  \item \textbf{损失正则性。} P2 使用 \(\ell\) 在相关残差范围内的 Lipschitz 界；平方损失不是全局 Lipschitz，因此 P2 对平方损失只作局部解释，P3 给出精确平方损失计算。
  \item \textbf{噪声模型。} P4 的线性无偏估计界假设零均值、协方差为 \(\eta^2 I\) 的噪声。若噪声相关，应改用 whitened operator 或 Fisher information \(A^\top C_n^{-1}A\)。
  \item \textbf{不声明神经网络动力学。} 这些命题不证明 \eng{UNet} 或 \eng{SGD} 必然漂移，也不分析非凸训练动力学；它们只说明，一旦漂移分量位于 exact/near null space，有界观测域损失在该分量上提供零或趋零的信息与恢复力。
  \item \textbf{proxy 反相关是条件性的。} L5 要求扰动相对 \(u_0\) 更高频、边缘增强或至少增加分母能量，并且近似正交于固定 raw-control highpass map \(c\)。若条件不满足，一阶符号可为零或改变。
\end{itemize}

\subsection{FRC 方法学}

本节详述 Fourier Ring Correlation（\eng{FRC}）的实验方案与全部控制组结果。

\subsubsection{分层分半构造}

\eng{FRC} 的核心思路是将 248 帧 clean set 分为两个子集 A、B，各自独立重建后计算频域相关。
为保证分半不引入系统性偏差，采用相位分层（phase-stratified）设计：将每帧的 stage command 经坐标变换映射为亚像素相位 \(\phi = \mathrm{mod}(\text{shift}, 1)\)，
按 \(5 \times 5 = 25\) 个 bin 分层（每 bin 7--13 帧），在每个 bin 内随机置换后奇偶交替分配给 A/B（约束 \(|N_A - N_B| \le 1\)）。
A 和 B 各自以 bilinear drizzle 重建到 \(5\times\) 诊断网格（hr\_pitch \(=\) 2\,\(\mu\)m），空 bin 填全局均值。
使用 seeds \{42, 123, 456\} 重复三次，主曲线取逐 ring nanmean。窗口函数为 Tukey \(\alpha = 0.25\)，边缘裁剪 16 \eng{LR} px。

\eng{FRC} 按环状积分计算：
\[
  \mathrm{FRC}(r) = \frac{\sum_{|\mathbf{f}| \in r} \mathrm{Re}\bigl(F_A(\mathbf{f}) \cdot \overline{F_B(\mathbf{f})}\bigr)}{\sqrt{\sum_{|\mathbf{f}| \in r} |F_A(\mathbf{f})|^2 \cdot \sum_{|\mathbf{f}| \in r} |F_B(\mathbf{f})|^2}}.
\]
cutoff 判据采用 \(1/7\) 常数阈值（0.142857）和 half-bit 曲线，取第一个 \eng{FRC} 低于阈值的 ring 对应周期。

\subsubsection{主结果与控制组}

\(1/7\) cutoff 为 17.03\,\(\mu\)m（3-seed 均值曲线），逐 seed 分别为 16.17/16.17/17.03\,\(\mu\)m（std \(=\) 0.50\,\(\mu\)m）；half-bit 判据给出相同值。
本文仅主张 17.0\,\(\mu\)m：超过 20\,\(\mu\)m 分辨率的相干信息确实存在，但低于 11--14\,\(\mu\)m 的理论期望。

四个控制组的设计与结果如下。正控制（单帧 bicubic 加噪）的 cutoff 为 13.58\,\(\mu\)m，未达到「明显差于 main」的预期，
原因是单帧插值的平滑频谱在低噪声下自相关偏高；负控制（shift-shuffle）在 8--12\,\(\mu\)m 段中位 \eng{FRC} 仍有 0.504，
因为置换后仍共享场景低频与网格结构，部分失效。漂移控制（前/后半按采集序分别重建）cutoff 退化至 26.20\,\(\mu\)m，
符合预期——时间间隔放大了热漂移的影响。zero-coverage 统计显示空 bin 均值 27.2\%、最大 36.2\%，表明 \(5\times\) 网格存在欠覆盖。

\subsubsection{10--12 \texorpdfstring{\(\mu\)}{u}m 反弹}

在 10--12\,\(\mu\)m 周期段，\eng{FRC} 出现异常高值，需要谨慎判读。关键频带的对比数据如下。

\begin{table}[ht]
  \centering
  \caption{关键周期段的 \eng{FRC}：main 与三个控制组对比}
  \label{tab:frc-band}
  \small
  \begin{tabular}{rcccc}
    \toprule
    周期 (\(\mu\)m) & main & bicubic 正控 & shuffle 负控 & drift 控 \\
    \midrule
    16 & 0.138 & 0.691 & \(-0.012\) & 0.117 \\
    12 & 0.593 & 0.002 & \(-0.095\) & 0.578 \\
    10 & 0.935 & 0.012 & 0.906 & 0.887 \\
    8  & 0.545 & \(-0.023\) & 0.390 & 0.345 \\
    \bottomrule
  \end{tabular}
\end{table}

10--12\,\(\mu\)m 的高 \eng{FRC} 同时出现在负控制（10\,\(\mu\)m 处 0.906）和漂移控制（0.887）中，
说明该反弹由 coverage/lattice 结构与漂移驱动，而非真实分辨率信息。因此不作分辨率证据使用。

\subsubsection{MAP-TV 前后对照}

以 \eng{MAP-TV}（\(\sigma = 0.2\), \(\lambda = 10^{-3}\)）对 drizzle 输出做去卷积后，split-half \eng{FRC} 在 20--10\,\(\mu\)m 段从 0.319/0.088/0.053/0.575/0.893 提升至 0.976/0.965/0.955/0.947/0.934。
该提升反映的是分半一致性（稳定性）的改善，而非光学分辨率本身的提升。

\subsection{标定不确定度传播}

\subsubsection{旋转角 \texorpdfstring{\(\theta\)}{theta} 的误差传播}

坐标到位移的映射为
\[
  d_x = \frac{X\cos\theta + Y\sin\theta}{10}, \quad d_y = \frac{-X\sin\theta + Y\cos\theta}{10}
\]
（单位 \(\mu\)m\(\rightarrow\)px）。对 \(\theta\) 求导，以最大步进 \(X = 40\)\,\(\mu\)m、\(\delta\theta = 0.1° = 1.745\)\,mrad 为例：
\(\delta d_x \approx 0.0052\)\,px，\(\delta d_y \approx 0.0047\)\,px，合成约 0.007\,px。
位移幅值 \(|s| = X/10 = 4.0\)\,px 与 \(\theta\) 一阶无关——\(\theta\) 误差表现为方向偏差而非步长误差。
该传播量比 alignment 残差（Chamfer 0.134\,px，见 B.2.2）低一个量级，不是误差预算的主项。

\eng{AVI} 的独立方向验证（gradient-NCC 合并估计 \(\theta = 47.14°\)，95\% CI [46.36°, 47.92°]）覆盖了标定值 47.6°，
但 X/Y 子组存在约 3° 系统差异（X 48.70° vs Y 45.63°），且 \eng{AVI} 为 8-bit 渲染视频，仅作一致性检验，不替换配置。

\subsubsection{PSF \texorpdfstring{\(\sigma\)}{sigma} 区间}

三条独立标定路线给出了不一致的结果：Route A（forward 残差）\(\sigma = 0.226\)\,px [0.208, 0.240]；
Route B（ESF 拟合）\(\sigma = 1.129\)\,px [1.041, 1.215]；Route C（joint hold-out）\(\sigma = 0.119\)\,px。
三者 spread 达 1.01\,px，远超 0.05\,px 容差。

\eng{EP15 M3} 仲裁澄清了分歧的物理根源：Route B 的 1.129\,px 实际上是 \eng{PSF} 与热/几何边缘宽度的卷积表观值，满足 \(\sigma_{\mathrm{total}}^2 = \sigma_{\mathrm{PSF}}^2 + w_{\mathrm{edge}}^2\)。
这意味着边缘表观宽度不能直接等同于 \eng{PSF}，同时也为反卷积的激进度设定了物理上限——即使完全去除光学 \eng{PSF}，热扩散造成的边缘宽度仍然存在。
综合 \eng{FRC} 形状拟合（\(\sigma = 0.2\) 最优）与多边缘 ESF 上界（0.546\,px），最终采纳区间为 \(\sigma \in [0.2, 0.5]\) \eng{LR} px。

该区间的传播后果显著：\(2\times\) Nyquist 处 \eng{MTF} 在区间内变化 65 倍（0.454\(\rightarrow\)0.007），因此所有依赖 \eng{PSF} 的计算均须把 \(\sigma\) 作为扫描参数。

\subsubsection{误差预算汇总}

\begin{table}[ht]
  \centering
  \caption{误差预算汇总}
  \label{tab:error-budget}
  \small
  \begin{tabular}{p{0.24\linewidth}p{0.32\linewidth}p{0.34\linewidth}}
    \toprule
    误差源 & 量级 & 对结论的影响 \\
    \midrule
    \(\theta\) \(\pm\)0.1° & \(\sim\)0.007\,px @40\,\(\mu\)m & 可忽略（不足 alignment 残差的 5\%） \\
    alignment 残差 & Chamfer 0.134\,px（refined） & 主要几何误差项；E3 消融显示改善端到端 corr \(+\)0.11（见 D.5.3） \\
    PSF \(\sigma\) 区间 & [0.2, 0.5]\,px \(\rightarrow\) \(2\times\) MTF 65\(\times\) 跨度 & 反卷积与合成均按区间扫描 \\
    噪声底 & 0.0724\,\(^\circ\)C & 可恢复对比度的下限 \\
    热漂移 & session 内 \(-0.60\)\,\(^\circ\)C；跨 session 中位 2.91\,\(^\circ\)C & session 门控 \(+\) 分层 split 的设计动机（见 B.1、B.3） \\
    \bottomrule
  \end{tabular}
\end{table}

\input{supp_sections/B_system_data}
\section{方法实现细节}

\subsection{TCForge 物理匹配合成平台}

本节详述主文~\S4.2 中合成数据平台 \eng{TCForge} 的实现，其目标是生成与真实 \eng{LWIR} 微扫描在物理特性上匹配的训练数据。

\subsubsection{场景几何生成}

\eng{TCForge} 通过 8 层随机原语叠加构造 \eng{IC} 布局的几何 mask：依次放置大块、中宽柱、L 形走线（25\% 粗 / 75\% 细）、框结构、细部与引脚阵列、圆形 via、对边 pin，最后做减通道处理。
所有层均由 seed 确定性控制。特征尺度由难度参数决定：easy 的 major/minor/line 为 80/60/45\,\(\mu\)m，medium 为 55/42/30\,\(\mu\)m，hard 为 35/28/22\,\(\mu\)m，stress 为 22/16/12\,\(\mu\)m。
场景整体旋转角为 \(\theta = 47.6° + U(-1.5°, +1.5°)\)，围绕实测 stage 角度抖动以匹配真实几何取向。

为消除阶梯锯齿对训练的污染，默认启用 4\(\times\)SSAA 抗锯齿：在 4 倍画布上绘制后旋转，再做 \(4 \times 4\) 块均值降采，得到 \(\text{coverage} \in [0, 1]\) 的连续值。
这样 \eng{HR} 目标不含阶梯锯齿，亚像素细线携带 coverage 缩放后的幅度。

\subsubsection{温度渲染与物理随机化}

温度场渲染公式为
\[
  T = T_{\mathrm{bg}} + \Delta T \cdot \text{coverage} + A \cdot \text{smooth\_noise}.
\]
关键物理参数的随机化分布经过仔细设计以覆盖实测范围：\(\Delta T \sim U(0.5, 5.0)\;^\circ\mathrm{C}\)（覆盖实测内/外轮廓对比度 1.94/2.49\,\(^\circ\)C）；
\eng{PSF} \(\sigma \sim U(0.15, 0.55)\) \eng{LR} px（覆盖采纳区间 [0.2, 0.5]，见 A.5.2），其中 30\% 使用椭圆 \eng{PSF}、10\% 使用 Airy 函数替代 Gaussian；
噪声遵循 lognormal 分布，均值等于实测噪声底 0.0724\,\(^\circ\)C；漂移幅度 \(\sim U(0, 0.3)\;^\circ\mathrm{C}\)；
亚像素位移使用实测 contour\_refined 的 248 帧 profile 重放并叠加 \(\sigma = 0.02\)\,px 的 jitter。

Forward 退化模式默认采用 physical\_block\_average：对 \eng{HR} 图先做 \eng{PSF} blur，再对各 shift 执行分块双线性均值降采（scale 2 或 4）。

\subsubsection{训练池}

当前构建了两个训练池。training\_pool\_2x\_aa 包含 1000 个 scene（难度分布 easy 200/medium 400/hard 300/stress 100），每 scene 248 帧，
\eng{LR} \(480 \times 640\) \(\rightarrow\) \eng{HR} \(960 \times 1280\)，用于 \(1\times\) 统计输入的各训练变体：\eng{Stats}、\eng{Stats+PixelShuffle}、\eng{Stats+HP-FC}、\eng{Stats+Full-FC}（run ids: v8.1a/v8.1b/v9b/v9d）。
training\_pool\_2x\_aa\_burst 在此基础上保存了完整的 \eng{LR} burst（约 152\,GB），用于 hybrid 输入的各变体：\eng{Hybrid}、\eng{Hybrid+Native-FC}、\eng{Hybrid+ResObs}（run ids: V9A/V9C/V10）。

为提高训练效率，还预计算了 \(K = 4\) 组 drizzle 变体：\(k = 0\) 使用全部 248 帧且无 shift 噪声（canonical，与推理一致）；
\(k = 1 \ldots 3\) 保留 60--100\% 帧（\(\ge 30\)）并叠加 \(N(0, 0.05\;\text{px})\) 的 shift 噪声。
训练时按 \([\text{seed}, \text{epoch}, \text{scene}, \text{0xBEEF}]\) 确定性抽样变体编号，将帧子集与 shift 噪声增广烘焙入预计算结果中，避免训练期间现场 drizzle 阻塞 dataloader。

\subsection{网络与损失}

\subsubsection{网络架构}

\eng{ThermalSRUNet} 采用编码-解码结构：编码路径为 3 层 MaxPool \(+\) ConvBlock（每块 \(2 \times 3 \times 3\) Conv \(+\) GroupNorm \(+\) SiLU \(+\) SE attention），
通道数按 \(c, 2c, 4c, 8c\) 递增（\(c = 64\)）；解码路径使用双线性上采样 \(+\) skip concat \(+\) ConvBlock。
\eng{HR} 输出头默认采用 bilinear 模式（双线性 \(\times\) scale \(\rightarrow\) \(3 \times 3\) Conv），
\eng{Stats+PixelShuffle}（v8.1b）消融变体另测试了 pixelshuffle 模式（ICNR 初始化 sub-pixel conv \(+\) PixelShuffle \(+\) HRResBlock），
后者产生条纹伪影且 proxy 全程更差（见 D.4.1），仅作为 head 归因对照保留。

在 hybrid 输入模式下，网络在 \(2\times\) 网格上以等效 scale\,=\,1 运行，因为输入已经是 \(2\times\) 分辨率。

\subsubsection{损失函数}

\eng{ContourSRLoss} 的总损失为
\[
  \mathcal{L} = w_m \cdot \text{MSE} + w_{hp} \cdot \text{HP} + w_e \cdot \text{Edge} + w_s \cdot (1 - \text{SSIM}) + w_{gv} \cdot \text{GV} + w_{lap} \cdot \text{Lap} + w_{fm} \cdot \text{FM} + w_r \cdot R.
\]
其中各项含义如下。MSE 为标准均方误差，可乘 gap\_weight。HP 为 highpass 结构损失 \(\|HP(\hat{x}) - HP(x)\|_1\)，预处理 \(HP(x) = x - G_{\sigma=5}(x)\)，乘以 structure/thin/gap 三种权重图。
Edge 为梯度幅值损失 \(\||\nabla\hat{x}| - |\nabla x|\|_1\) 加 0.25 倍半分辨率版本。SSIM 使用 \(11 \times 11\) Gaussian 窗。
GV（grad-vector）为 \(\|g_{x,p} - g_{x,t}\|_1 + \|g_{y,p} - g_{y,t}\|_1\)，保留梯度的完整向量信息以捕捉边缘膨胀和畸变——这是幅值型 Edge loss 看不到的方向偏差。
Lap 为 \(\text{ReLU}(|\Delta x| - |\Delta\hat{x}|)\) 的均值，只罚「比目标钝」。
FM 为 forward consistency \(\text{MSE}(\text{AvgPool}_s(G_{\sigma s}(\hat{x})),\; y_{\text{obs}})\)，可选 full-band 或 highpass 模式。
R 为残差惩罚 \(w_r \cdot \text{mean}(|\hat{x} - x_{\text{input}}|)\)（\eng{Hybrid+ResObs} / V10 使用）。

权重图机制使损失在结构关键区域有针对性地加强：structure\_boost 按梯度强度对 HP/GV 权重图加权；thin\_boost 在宽度 \(\le 3\) \eng{HR} px 的细结构上倍增；gap\_boost 在两侧有结构的窄缝背景上倍增。

conservative 权重壳（自 \eng{Stats} / v8.1a 起冻结，FC 与 hybrid 系列沿用）设定为 \(w_m = 0.3\)、\(w_{hp} = 0.8\)、structure\_boost \(=\) 2.0、\(w_{gv} = 0.15\)、thin\_boost \(=\) 3.0、gap\_boost \(=\) 2.0。
\eng{HotLoss}（v6）使用更激进的 hot 壳（\(w_{hp} = 1.0\)、structure\_boost \(=\) 4.0 等），叠加 Lap 0.1 和 full-band FM 0.1，保留为漂移放大参照变体。

\subsubsection{Hybrid 8 通道输入}

\eng{Hybrid}、\eng{Hybrid+Native-FC}、\eng{Hybrid+ResObs}（run ids: V9A/V9C/V10）使用 hybrid 8 通道输入，在 \(2\times\) 网格上融合 \(1\times\) 统计特征与 \(2\times\) drizzle 证据。
通道 0--4 为 \(1\times\) fused 统计（aligned mean/median/coverage/variance/highpass fused）经双线性上采到 \(2\times\)；
通道 5 为 drizzle mean @\(2\times\)，是观测域证据的主通道（也是 \eng{Hybrid+ResObs} 残差惩罚的基准）；
通道 6 为 drizzle coverage @\(2\times\)（空洞可识别：未观测 bin 以全局均值填充，coverage \(=\) 0）；通道 7 为 drizzle variance @\(2\times\)。

\eng{Hybrid+Native-FC}（V9C）的合法锚配套值得说明：hybrid 的通道 0 是上采样均值而非合法的 \(1\times\) 观测，
因此 forward 锚需另走数据管线携带原生 \(1\times\) aligned-mean patch（偶数 \(2\times\) origin 裁剪、增广同步）。

\subsection{训练配置对照}

七个变体的训练配置共享以下设定：scale\,=\,2、base\_channels\,=\,64、patch\_size\_hr\,=\,256、total\_steps\,=\,60000、lr\,=\,\(2 \times 10^{-4}\)、AMP、compile、
edge/ssim/coarse 权重 0.05/0.15/0.25、highpass\_sigma\,=\,5、real\_eval 使用 248 帧 contour\_refined \(+\) center-1/3 \(+\) zoom3。

各变体的差异字段构成一个 \(\{\text{input}\} \times \{\text{anchor}\}\) 消融矩阵。

\begin{table}[ht]
  \centering
  \caption{七个变体的差异字段对照}
  \label{tab:variant-config}
  \scriptsize
  \resizebox{\linewidth}{!}{%
  \begin{tabular}{lccccccc}
    \toprule
    字段 & HotLoss & Stats & Stats+PixelShuffle & Stats+HP-FC & Stats+Full-FC & Hybrid & Hybrid+Native-FC \\
    \midrule
    run id      & v6        & v8.1a    & v8.1b        & v9b      & v9d      & V9A             & V9C \\
    训练池      & 2x        & 2x\_aa   & 2x\_aa       & 2x\_aa   & 2x\_aa   & 2x\_aa\_burst   & 2x\_aa\_burst \\
    input / in\_ch & lr / 5 & lr / 5   & lr / 5       & lr / 5   & lr / 5   & hybrid / 8      & hybrid / 8 \\
    HR 头       & bilinear  & bilinear & pixelshuffle & bilinear & bilinear & bilinear        & bilinear \\
    loss 壳     & hot+Lap   & cons.    & cons.        & cons.    & cons.    & cons.           & cons. \\
    forward 锚  & full 0.1  & 无       & 无           & hp 0.1   & full 0.1 & 无              & legal-1x hp 0.1 \\
    batch\_size & 128       & 128      & 128          & 128      & 128      & 64              & 64 \\
    \bottomrule
  \end{tabular}
  }
\end{table}

消融矩阵的直读方式：\(\{1\times, \text{hybrid}\} \times \{\text{none}, \text{HP-FC}, \text{Full-FC}, \text{Native-FC}\}\) \(=\) \eng{Stats} / \eng{Stats+HP-FC} / \eng{Stats+Full-FC} / \eng{Hybrid} / \eng{Hybrid+Native-FC}（\eng{HotLoss} 和 \eng{Stats+PixelShuffle} 分别为 loss 温度和 head 的额外归因变体）。

\textbf{可比性说明：} \eng{Hybrid}（V9A）前 35K 步使用 batch\_size\,=\,128，之后与 \eng{Hybrid+Native-FC}（V9C）全程 batch\_size\,=\,64——\eng{Hybrid} 与 \eng{Hybrid+Native-FC} 两个变体与 \(1\times\) 各变体的训练动力学不完全同条件。
\eng{Hybrid}（V9A）30K 保真悬崖与 batch size 切换时间重合，混杂未排除。

\subsection{经典方法}

\subsubsection{Drizzle}

Drizzle 使用 STScI drizzle 库实现，pixfrac 默认 0.7（sweep \{1.0, 0.8, 0.7, 0.6, 0.5\}），kernel 为 square，输出网格 \((H \times 2, W \times 2)\)。
coverage \(< 1.0\) 的像素置 NaN。\eng{TCForge} 训练侧的 scatter 使用 bilinear kernel，两种 kernel 的差异在 D.7 的 fine-window 指标中以 drizzle 输入通道参照点的形式体现。

\subsubsection{各向异性 TGV}

\eng{TGV} 重建采用外层 FISTA 框架 \(x_{k+1} = \text{TGV\_prox}(z_k - \eta \nabla D(z_k))\)，针对本数据集的 raster 各向异性（见 B.3.1）做了两项关键设计。

\textbf{第一，各向异性对偶投影。} \eng{TGV} 内层（Chambolle-Pock 路径）的一阶对偶球改为椭圆 \(\{(a,b) : (a/r_a)^2 + (b/r_b)^2 \le 1\}\)，
其中 \(r_a = \alpha_1 \cdot r_y\)、\(r_b = \alpha_1\)（\(r_y = 1.5\)）；二阶对称张量同构造。这使得 Y 方向正则强度弱于 X 方向，补偿 Y 方向数据约束的不足。

\textbf{第二，coverage 加权数据梯度。} 每个 \eng{HR} 像素的数据保真梯度除以预计算的 coverage（由 bilinear splat \(+\) \eng{PSF} adjoint 得到），而非统一除以帧数 \(N\)。
这修正了 bilinear scatter 把权重集中在固定 \eng{HR} 行上所导致的水平条纹。

参数网格为 \(\lambda_{\text{tv}} \in \{3 \times 10^{-4}, 10^{-3}, 3 \times 10^{-3}\} \times \sigma_{\text{PSF}} \in \{0.18, 0.50\}\)，\(\alpha_{\text{ratio}} = 2.0\)（\(\alpha_1 = \lambda\), \(\alpha_0 = \lambda/2\)），外层 100 iter、内层 80 iter。
最优参数下 artifact 从 3.870 降至 0.695（\(-82\%\)），raw-control corr 从 0.902 升至 0.916，耗时 30.8\,min CPU。

\subsubsection{MAP-TV 去卷积基准}

\eng{MAP-TV} 基准用于验收门控（C.5.0 采纳规则）：学习方法必须在 \eng{FRC} 频带一致性与 zigzag profile 指标上同时不输 \eng{MAP-TV} 才可采纳。
其 forward 模型为 shift \(\rightarrow\) Gaussian \eng{PSF}（\(\sigma\) \eng{LR} px）\(\rightarrow\) avg\_pool box，使用 GPU FISTA \(+\) smoothed-TV 梯度优化，150 iter。

参数网格为 \(\sigma \in \{0.2, 0.3, 0.4, 0.5\} \times \lambda \in \{3 \times 10^{-4}, 10^{-3}, 3 \times 10^{-3}\}\)。
选择规则为每个 \(\sigma\) 内按 split-half proxy（split\_half\_nrmse \(+\) 0.05\(\times\)artifact \(+\) 0.08\(\times\)std\_excess）取最小，再跨 \(\sigma\) 选全局最优，结果为 \(\sigma = 0.2\)、\(\lambda = 10^{-3}\)，耗时 4563\,s GPU。

\subsection{Checkpoint 选择协议}

训练期间的 checkpoint 选择对学习变体的最终输出质量至关重要。本节详述主文~\S4.5 和~\S5.6 的选择规则与采纳判定。

\subsubsection{采纳规则（Adoption rule）}

一个学习方法被采纳，当且仅当在其选定 checkpoint 上同时满足：
（a）\eng{FRC} 频带一致性 \(\ge\) \eng{MAP-TV} 基准；（b）zigzag FWHM/dip 不差于基准；（c）proxy 轨迹无 post-selection 漂移悬崖；（d）双域视觉 panel 均通过。
否则经典基准仍是交付物。正是这条保守规则，让诚实的负结果（PixelShuffle 头、\(4\times\) 网络、loss 侧锚定）与阳性结果得以并列报告——拒绝一个方法和采纳一个方法用的是同一把尺子。

\subsubsection{选择流程}

对每个变体的所有 eval checkpoint，首先归一化两个 proxy：\(a_{\text{norm}} = (\text{artifact} - \min) / (\max - \min)\)，\(c_{\text{norm}} = (\max_{\text{corr}} - \text{corr}) / (\max_{\text{corr}} - \min_{\text{corr}})\)（corr 反向）。
计算到理想点 \((0, 0)\) 的距离 \(d = \sqrt{a_{\text{norm}}^2 + c_{\text{norm}}^2}\)，按 \(d\) 升序遍历。
为避免选中相邻 checkpoint 的冗余候选，采用 5K 窗口去重：每个 checkpoint 归属 window \(=\) step // 5000，同窗只保留最先入选者，最多保留 3 个候选。
60K 步若未入选则强制追加为漂移参照。Rank-1 为 canonical 候选。最终每个候选需通过温度域三联 panel 的视觉 gate，人工过门后定稿。

\subsubsection{已执行选点}

\eng{EP11} 四个变体的选点结果如下：\eng{HotLoss}（v6）canonical 8K（artifact/corr \(=\) 0.330/0.774），\eng{Stats}（v8.1a）canonical 15K（0.392/0.758），
\eng{Stats+PixelShuffle}（v8.1b）canonical 5K（0.370/0.739），\eng{Stats+HP-FC}（v9b）canonical 11K（0.339/0.777）。\eng{TGV} Pareto 参考点为 (0.695, 0.916)。
各变体 60K 端点的 artifact 均显著恶化——按端点上报会让每个变体交出最差 checkpoint，这正是选择协议存在的意义。

统一 harness 后的补充选点为：\eng{Hybrid}（V9A）canonical 10K（harness artifact/corr \(=\) 1.762/0.719），\eng{Hybrid+Native-FC}（V9C）canonical 5K（1.669/0.718），
\eng{Stats+Full-FC}（V9D）canonical 7K（1.726/0.771），\eng{Hybrid+ResObs}（V10）采用高-\(\lambda\) sweep 的工作点 \(\lambda=1.2\)@15K（2.726/0.711；fine-window 为 hp\_corr\_input 0.922、sharp\_p95 0.987、lattice 0.0141）。
\eng{Hybrid}（V9A）60K 不作为交付 checkpoint，仅作为 F5 late-drift visual control。
需强调 hybrid 变体与 \(1\times\) 变体的 proxy 不可跨列横比（见 A.3.2）；最终 T1/T2 横评只使用 \texttt{output/ep11\_unified\_harness/} 的同一 harness scale。

\subsection{结构级指标与视觉门控}

本节详述主文实验中使用的结构级评估指标与视觉门控流程。

\subsubsection{Zigzag 走线 profile 指标}

以客户关心的中心走线为测量对象。在固定横截面上，对每条 profile 测表观 FWHM（full width at half maximum）与谷深（dip），报告 median FWHM/dip 及 per-profile 分布。
当 per-profile 结果混杂时如实报告为「混杂」，不取巧。需要注意，\(2\times\) 网格下量化步长为 5\,\(\mu\)m，FWHM 可能饱和（主文~\S5.1 多个方法同为 40\,\(\mu\)m），区分度低时不应过度解读。

\subsubsection{双域视觉 panel}

双域 panel 由两张并列图组成。第一张为高通结构图（边缘证据），使用 \(HP(x) = x - G_{\sigma=5}(x)\)，红/蓝分别表示相对局部背景的正/负有符号响应，白色表示无变化。
第二张为普通温度图（inferno colormap，1--99 百分位裁剪），专门用于捕捉「只有边缘变亮但温度场整体失真」这类被高通图美化掉的失效模式。
两图并列能让审查者同时看到高频结构证据和低频温度合理性。

\subsubsection{对齐门控一致性}

重建不得劣化被门控锚点上的 hold-out 轮廓一致性。该角色由 \eng{EP04} 对齐质量评估承担：
若某方法在门控锚点上的 split-half 轮廓一致性低于对照，该方法在该区域不予采纳。

\section{完整实验结果}

\subsection{统一 benchmark 全表}

统一 harness 已于 2026-06-14 完成，入口脚本为 \texttt{algos/ep11\_dl\_benchmark/scripts/run\_unified\_harness\_t1\_t2.py}，产物位于 \texttt{output/ep11\_unified\_harness/}。
完整列（runtime、source、cache path、temperature sanity、FRC@20/10 等）保存在 \texttt{all\_arm\_metrics.csv}；主文 T1/T2 分别使用 \texttt{t1\_metrics.csv} 与 \texttt{t2\_metrics.csv}。
所有 artifact 数字均为 EP11/common.metrics harness scale，不能与 TensorBoard \texttt{eval\_real/*} artifact 混表；
\texttt{tb\_vs\_harness\_scale\_check.csv} 给出三条对照（如 \eng{Stats+HP-FC} / v9b@11K: TB 0.3385 vs harness 1.7662）。

\subsubsection{T1 主表精选列}

\begin{table}[ht]
  \centering
  \caption{T1 主表精选列（统一 harness）}
  \label{tab:d-t1}
  \scriptsize
  \resizebox{\linewidth}{!}{%
  \begin{tabular}{lrrrrrrrrrr}
    \toprule
    方法 & step & split\(\downarrow\) & artifact\(\downarrow\) & corr\(\uparrow\) & FRC16\(\uparrow\) & FRC14\(\uparrow\) & FRC12\(\uparrow\) & lattice\(\downarrow\) & sharp\(\uparrow\) & FWHM\,\(\mu\)m\(\downarrow\) \\
    \midrule
    bicubic                          & ---  & 0.047 & 1.388 & 1.000 & 0.951 & 0.945 & 0.941 & 0.001 & 0.344 & 60 \\
    drizzle                          & ---  & 0.023 & 1.138 & 0.771 & 0.439 & 0.564 & 0.765 & 0.001 & 0.481 & 45 \\
    MAP-TV (\(5\times\))             & ---  & 0.024 & 2.333 & 0.438 & 0.965 & 0.955 & 0.947 & 0.002 & 0.352 & 42 \\
    TGV                              & ---  & 0.032 & 0.695 & 0.741 & 0.975 & 0.969 & 0.955 & 0.011 & 0.999 & 40 \\
    HotLoss                          & 8K   & 0.051 & 1.789 & 0.774 & 0.752 & 0.761 & 0.217 & 0.001 & 0.656 & 40 \\
    Stats                            & 15K  & 0.073 & 1.943 & 0.758 & 0.711 & 0.660 & 0.429 & 0.001 & 0.864 & 40 \\
    Stats+HP-FC                      & 11K  & 0.054 & 1.766 & 0.777 & 0.744 & 0.741 & 0.197 & 0.001 & 0.698 & 40 \\
    Stats+Full-FC                    & 7K   & 0.054 & 1.726 & 0.771 & 0.744 & 0.708 & 0.150 & 0.001 & 0.697 & 40 \\
    Hybrid                           & 10K  & 0.054 & 1.762 & 0.719 & 0.945 & 0.870 & 0.190 & 0.001 & 0.683 & 40 \\
    Hybrid+Native-FC                 & 5K   & 0.064 & 1.669 & 0.718 & 0.891 & 0.769 & 0.195 & 0.001 & 0.766 & 40 \\
    Hybrid+ResObs (\(\lambda=1.2\))  & 15K  & 0.041 & 2.726 & 0.711 & 0.986 & 0.984 & 0.979 & 0.013 & 0.968 & 40 \\
    \bottomrule
  \end{tabular}
  }
\end{table}

边界说明：\eng{MAP-TV} 是预计算 \(5\times\) anchor；\eng{TGV} 的 split/FRC 列已用独立 \eng{TGV} half-set 重建回填
（\texttt{output/ep11\_unified\_harness/tgv\_split\_frc.json}，seed 42；full-run self-check 相对 L2 \(=\) 0.0），不再是 \eng{EP16} drizzle proxy。
该 \eng{FRC} 仍是 split-consistency proxy，不是物理分辨率证据。
\eng{Hybrid+ResObs} / \eng{Hybrid} / \eng{Hybrid+Native-FC}（run ids: V10/V9A/V9C）的缓存均值约 23\,\(^\circ\)C，说明 hybrid / ResObs 推理没有漏加 drizzle base。

\subsubsection{T2 input × anchor 精选列}

\begin{table}[ht]
  \centering
  \caption{T2 输入 \(\times\) 锚定精选列}
  \label{tab:d-t2}
  \scriptsize
  \resizebox{\linewidth}{!}{%
  \begin{tabular}{llrrrrrrr}
    \toprule
    方法 & input & anchor/参数化 & step & split\(\downarrow\) & artifact\(\downarrow\) & corr\(\uparrow\) & FRC16\(\uparrow\) & FRC14\(\uparrow\) \\
    \midrule
    Stats              & 1x stats          & none                 & 15K & 0.073 & 1.943 & 0.758 & 0.711 & 0.660 \\
    Stats+PixelShuffle & 1x stats          & PixelShuffle control & 5K  & 0.050 & 1.782 & 0.739 & 0.734 & 0.692 \\
    Stats+HP-FC        & 1x stats          & HP-FC 0.1            & 11K & 0.054 & 1.766 & 0.777 & 0.744 & 0.741 \\
    Stats+Full-FC      & 1x stats          & Full-FC 0.1          & 7K  & 0.054 & 1.726 & 0.771 & 0.744 & 0.708 \\
    Hybrid             & hybrid drizzle2x  & none                 & 10K & 0.054 & 1.762 & 0.719 & 0.945 & 0.870 \\
    Hybrid+Native-FC   & hybrid drizzle2x  & Native-FC 0.1        & 5K  & 0.064 & 1.669 & 0.718 & 0.891 & 0.769 \\
    Hybrid+ResObs      & hybrid drizzle2x  & ResObs \(\lambda=1.2\) & 15K & 0.041 & 2.726 & 0.711 & 0.986 & 0.984 \\
    \bottomrule
  \end{tabular}
  }
\end{table}

\noindent（FRC12 与 sharp 列见 \texttt{t2\_metrics.csv}；Hybrid+ResObs FRC12 \(=\) 0.979、sharp \(=\) 0.968，其余变体 sharp 见 T1。）

读数：输入端引入 drizzle2x 证据能改变 fine-window 结构权衡，但不是自动保真胜利；loss-side anchor 在 \(1\times\)/full/band/legal-hybrid 变体中均未压平真实数据漂移。
\eng{Hybrid+ResObs}（V10）把 sharp/FRC 推高，同时 artifact 与 corr 变差；它是可调 trade-off，不是可认证赢家。

\subsection{FRC 完整档案}

本节归档 \eng{FRC} 分析的全部数值结果，方法学细节见 A.4。

\subsubsection{主曲线与 cutoff}

\(1/7\) cutoff 为 17.03\,\(\mu\)m（3-seed 均值曲线），逐 seed 分别为 16.17/16.17/17.03\,\(\mu\)m（std \(=\) 0.50\,\(\mu\)m），half-bit 判据给出相同值。
控制组 cutoff 为：bicubic 正控 13.58\,\(\mu\)m（未通过预期）、漂移控制 26.20\,\(\mu\)m；zero-coverage 均值 27.2\%、最大 36.2\%。
频带全表与控制组详细分析见 A.4.3 和 A.4.4。

\subsubsection{MAP-TV 前后 split-half 对照}

以 \eng{MAP-TV}（\(\sigma = 0.2\), \(\lambda = 10^{-3}\)）对 bare drizzle 做去卷积后，split-half \eng{FRC} 在 20/16/14/12/10\,\(\mu\)m 周期段分别从 0.319/0.088/0.053/0.575/0.893 提升至 0.976/0.965/0.955/0.947/0.934。
这反映的是 split-half 一致性的改善，而非光学分辨率本身的提升——bare cutoff 仍为 17.03\,\(\mu\)m，与 M2 一致。

\subsubsection{Zigzag 轮廓指标}

\eng{MAP-TV} 基准在 zigzag 剖面上的结构收益有限：中位 FWHM 从 114\,\(\mu\)m 缩窄至 100\,\(\mu\)m，dip 从 0.929 微升至 0.934，3/3 剖面保持分离但改善不均（2 条变宽、1 条显著变窄）。
论文措辞限定为 "limited contour enhancement"。

\subsection{漂移演化档案}

\subsubsection{全局漂移指标}

下表记录了各变体在训练期间真实数据 eval（TB-scale，248 帧 contour\_refined）上的 artifact（\(\downarrow\)）与 raw\_control\_corr（\(\uparrow\)）轨迹（每格为 artifact / corr）。

\begin{table}[ht]
  \centering
  \caption{各变体真实数据 eval 漂移轨迹（artifact / corr）}
  \label{tab:drift-archive}
  \small
  \resizebox{\linewidth}{!}{%
  \begin{tabular}{lccccc}
    \toprule
    step & Stats & Hybrid & Stats+HP-FC & Stats+Full-FC & Hybrid+Native-FC \\
    \midrule
    10K & 0.390 / 0.756 & 0.446 / 0.719 & 0.369 / 0.758 & 0.379 / 0.758 & 0.516 / 0.714 \\
    20K & 0.476 / 0.729 & 0.514 / 0.702 & 0.486 / 0.735 & 0.575 / 0.642 & 0.689 / 0.655 \\
    30K & 0.602 / 0.703 & 0.660 / 0.663 & 0.611 / 0.709 & 0.615 / 0.694 & 0.686 / 0.675 \\
    40K & 0.627 / 0.698 & 0.656 / 0.665 & 0.640 / 0.697 & 0.672 / 0.681 & 0.688 / 0.672 \\
    60K & 0.643 / 0.689 & 0.646 / 0.669 & 0.655 / 0.688 & 0.677 / 0.677 & 0.695 / 0.669 \\
    \bottomrule
  \end{tabular}
  }
\end{table}

从轨迹形状看，\eng{Hybrid}（V9A）是唯一在 30K\(\rightarrow\)60K 段漂移趋平的变体（变化量 \(-0.014\)/\(+0.007\)，其余单调恶化）。
\eng{Stats+Full-FC}（V9D）比 \eng{Stats+HP-FC}（V9B）更差且 1K--28K 段震荡剧烈。
\eng{Hybrid+Native-FC}（V9C）在 30K 后没有恢复，60K 端点为 0.695/0.669，说明 hybrid 输入下的合法 \(1\times\) anchor 也不能压平漂移。
需要强调的是，hybrid 变体（\eng{Hybrid} / \eng{Hybrid+Native-FC}）与 \(1\times\) 变体的 proxy 不可跨列横比（见 A.3.2），此处只看各变体轨迹形状。

\subsubsection{Fine-window 训练时间轴}

为更精细地诊断训练动态，在 \(2\times\) 网格全幅 \(960 \times 1280\) 的中心细线窗口（rows 384:518, cols 478:674——中心两条细 zigzag 及周边粗线）上，跟踪四个指标：
\(\text{hp\_corr\_input}\)（highpass 域对 drizzle 输入通道的 Pearson，衡量保真性）、\(\text{hp\_corr\_tgv}\)（同口径对 \eng{TGV} 的 Pearson，交叉验证）、
\(\text{sharp\_p95}\)（温度图 P95 梯度幅值，衡量锐度）、\(\text{lattice\_score}\)（highpass 域 \(|f| > 0.35\) cyc/px 频段能量占比，捕捉格纹伪影）。

两个参照点为：drizzle 输入通道 \((1.000, 0.503)\)——观测域上限（模糊但零幻觉）；\eng{TGV} \((0.960, 0.959)\)——经典前沿工作点。

\eng{Hybrid}（V9A）的训练时间轴揭示了一个关键现象——30K 保真悬崖。\eng{Hybrid} 10K 的 \(\text{hp\_corr\_input}\) 为 0.970，20K 达峰 0.974，
此后断崖下跌：30K 降至 0.908 并在 40K--60K 焊死于 \(0.906 \pm 0.001\)。
同时，锐度 \(\text{sharp\_p95}\) 在 30K 后超过 \eng{TGV}（1.147 vs 0.959），并在 40K--60K 继续攀升至 1.21--1.25——锐度超过 \eng{TGV} 的区间恰好与去相关重合，指示过冲为幻觉驱动。
对照 \eng{Stats}（v8.1a）60K 同窗口为 0.926/0.936，\eng{Hybrid} 60K（0.925/0.935）与之无差别——hybrid 早期增益被训练后期完全抹平。

需注意 \eng{Hybrid}（V9A）的 35K 中断（batch size 从 128 切换至 64）与悬崖时间重合，混杂尚未排除（见 C.3）。

Fine-window 是局部诊断窗口，依赖 \eng{TGV} 参照，独立性低于 FRC/proxy 对——仅用于归因与选型，不替代全局评估协议。

\subsection{负结果档案}

本节完整记录四项负结果，为主文~\S5.7 提供数据背书。每项按「现象 \(\rightarrow\) 数字 \(\rightarrow\) 结论边界」结构组织。

\subsubsection{PixelShuffle HR 头}

\eng{Stats+PixelShuffle}（v8.1b）使用 pixelshuffle \eng{HR} 头替代默认的 bilinear 头（其余与 \eng{Stats} / v8.1a 相同），
出现了中等边框间条纹状亮色伪影，锯齿未改善，中心细线模糊程度与 bilinear 头相同。
TB-scale 数据显示 artifact 全程高于 \eng{Stats}（0.413\(\rightarrow\)0.709 vs 0.390\(\rightarrow\)0.643），corr 全程更低（0.747\(\rightarrow\)0.667 vs 0.758\(\rightarrow\)0.689）。

结论边界：head 归因失败——中心细线瓶颈不在解码头。该结果与 \eng{Stats}（v8.1a）的 loss-cooldown 对照共同构成主文~\S5.3 的两个对照归因：细线模糊对 loss 温度和 head 均不敏感，指向输入信息瓶颈。

\subsubsection{4x 网络}

\eng{EP12} \(4\times\) 网络在所有指标上均劣于 \eng{EP07} \(2\times\) 上采样方案：raw-control highpass Pearson 0.223 vs 0.389，artifact 0.535 vs 0.472。

该实验结果与 A.1.3 的理论预测一致：\(4\times\) Nyquist \eng{MTF} \(\le 0.042\)（\(\sigma = 0.2\)），\(\sigma \ge 0.35\) 时接近零；相位 bin 在 \(3\times\)/\(4\times\) 网格出现 collapse（见 B.2.3）。
\(4\times\) 无真实增益且理论上不应有——这是「先第一性原理后实验」工作流的展示案例。

\subsubsection{Forward 锚定（Stats+HP-FC / Stats+Full-FC / Hybrid+Native-FC）}

这些锚定变体的漂移曲线与无锚 \eng{Stats}（v8.1a）几乎重合。\eng{Stats+HP-FC}（V9B）的 forward loss 自 10K 步起贴底于 0.004--0.009，同期 artifact 持续上爬。
\eng{Stats+HP-FC} 40K\(\rightarrow\)60K 漂移 \(+0.0145\)/\(-0.0082\) vs \eng{Stats} \(+0.016\)/\(-0.009\)，几乎完全重合。
\eng{Stats+Full-FC}（V9D）60K（0.677/0.677）劣于 \eng{Stats+HP-FC}（0.655/0.688），且 1K--28K 段剧烈震荡——复现了全频低通梯度冲突的已知问题。

\eng{Hybrid+Native-FC}（V9C）是最后一条反对意见检验：在 hybrid 输入含 \(2\times\) drizzle 证据的条件下，forward anchor 不再偷用上采样 mean，而是通过数据管线携带合法的 \(1\times\) aligned-mean patch。
结果仍为 0.516/0.714 @10K \(\rightarrow\) 0.695/0.669 @60K，未压平后期漂移。

结论边界：无论频带选择和输入证据路径如何，\(1\times\) 观测域锚定对真实漂移不可见——机制即 A.2.3 Proposition 1。
\eng{Stats+Full-FC}（V9D）关闭了「band 太窄」的反对意见；\eng{Hybrid+Native-FC}（V9C）关闭了「hybrid 第 0 通道不是合法 \(1\times\) 观测」的反对意见。
loss-side anchoring route 在 HP-FC / Full-FC / Native-FC 三种变体下正式关闭。

\subsubsection{AVI 作 SR 输入}

8-bit 渲染视频、约 67\% 重复帧、无温度矩阵（见 B.1.3）——直接排除。\eng{AVI} 仅用于方向一致性检验。
该负结果的意义在于：「能拿到的数据」不等于「能进重建的数据」，审计应先于算法。

\subsection{帧预算与鲁棒性}

本节报告 \eng{EP16} 的三组经典方法实验：E1 帧数预算、E2 shift 扰动、E3 对齐源消融。共 37 个 unique run 全部成功。
注意这是推理期稳定性研究，不属于统一 harness。该轨迹图原为主文 F7，现移至此处仅作 supp 参照。

\begin{figure}[ht]
  \centering
  \includegraphics[width=\linewidth]{../../output/paper_figures/fig07_budget_robustness.png}
  \caption{帧数预算（左）与 shift 扰动鲁棒性（右）。\eng{Drizzle} 的 raw-control corr 多数增益在 \(N = 62\) 即出现，
  FRC/split-half proxy 随更多相位覆盖继续改善；\eng{TGV} 全程 artifact 低于 drizzle 但 corr 非单调。
  横轴对应 \(N \in \{31,62,124,248\}\) 与 \(\sigma \in \{0,0.05,0.1,0.2\}\)\,px，详见下表。}
  \label{fig:d-budget}
\end{figure}

\subsubsection{帧数预算}

帧数预算实验使用 phase-stratified 子集（\(N \in \{31, 62, 124, 248\}\)，seeds 101/202/303），覆盖 drizzle 和 \eng{TGV} 两种经典方法。

\begin{table}[ht]
  \centering
  \caption{帧数预算扫描（drizzle 与 TGV）}
  \label{tab:budget}
  \small
  \resizebox{\linewidth}{!}{%
  \begin{tabular}{lrrrrrr}
    \toprule
    方法 & \(N\) & corr & split-half NRMSE & artifact & FRC@16\,\(\mu\)m & zigzag FWHM (\(\mu\)m) \\
    \midrule
    drizzle & 31  & 0.747\(\pm\)0.032 & 0.0715\(\pm\)0.0012 & 1.649\(\pm\)0.041 & 0.109\(\pm\)0.065 & 56.7\(\pm\)20.2 \\
    drizzle & 62  & 0.772\(\pm\)0.016 & 0.0514\(\pm\)0.0034 & 1.536\(\pm\)0.044 & 0.248\(\pm\)0.070 & 56.7\(\pm\)20.2 \\
    drizzle & 124 & 0.770\(\pm\)0.010 & 0.0361\(\pm\)0.0010 & 1.341\(\pm\)0.038 & 0.332\(\pm\)0.018 & 45.0 \\
    drizzle & 248 & 0.771 & 0.0306 & 1.145 & 0.479 & 45.0 \\
    TGV & 31  & 0.728\(\pm\)0.053 & --- & 0.946\(\pm\)0.003 & --- & 42.5\(\pm\)3.5 \\
    TGV & 62  & 0.754\(\pm\)0.009 & --- & 0.851\(\pm\)0.051 & --- & 40.0 \\
    TGV & 124 & 0.735\(\pm\)0.013 & --- & 0.747\(\pm\)0.011 & --- & 40.0 \\
    TGV & 248 & 0.741 & --- & 0.708 & --- & 40.0 \\
    \bottomrule
  \end{tabular}
  }
\end{table}

\eng{TGV} 的 split-half/FRC 列使用同子集 drizzle proxy（预算考虑），以 "---" 标注。

Drizzle corr 的增益大半在 \(N = 62\) 前到位（0.747\(\rightarrow\)0.772），之后趋平；split-half、artifact 和 \eng{FRC} 随 \(N\) 单调改善，受益于相位覆盖的累积。
\eng{TGV} artifact 全程低于 drizzle（0.946\(\rightarrow\)0.708），但 corr 呈非单调变化（0.728\(\rightarrow\)0.754\(\rightarrow\)0.735\(\rightarrow\)0.741）。

\subsubsection{Shift 扰动}

扰动实验在 \(\sigma \in \{0, 0.05, 0.1, 0.2\}\)\,px（seeds 401--403）下测试对齐误差的影响。
Drizzle corr 从 0.771 到 0.770 几乎不变，而 artifact 从 1.145 升至 1.434、FRC@16 从 0.479 降至 0.340。\eng{TGV} corr 同样稳定（0.741\(\rightarrow\)0.744）。

结论：鲁棒性是 metric-specific 的——raw-control corr 对 \(\le 0.2\)\,px 扰动几乎不敏感，而 coverage/FRC 类指标敏感。这是压力测试口径，非真实对齐误差估计。

\subsubsection{对齐源消融}

将 command\_prior 与 contour\_refined 两种对齐源的端到端效果作比：drizzle corr 从 0.662 升至 0.771（\(+0.109\)），FRC@16 从 0.0166 升至 0.479；\eng{TGV} corr 从 0.642 升至 0.741（\(+0.099\)）。

数据驱动对齐对两个经典方法都带来约 \(+0.10\) corr 的端到端增益，FRC@16 从几乎为零到 0.479——这是 B.2 对齐链投入的最终回报证据。

\subsection{视觉 gate panel}

已有 \eng{EP11} 四个变体（\eng{HotLoss}、\eng{Stats}、\eng{Stats+PixelShuffle}、\eng{Stats+HP-FC}）的 checkpoint 选择 panel：每个变体 canonical \(+\) 60K 对照的温度域三联图，标注 proxy 值。
\eng{Stats+PixelShuffle} 行保留条纹伪影证据。主文 F5 已补 \eng{Hybrid} late endpoint 60K 与 \eng{Hybrid+ResObs} \(\lambda=1.2\)@15K 的双域视觉 gate，
并把中心 \eng{ROI} 与 held-out ROI2 合成 4 行图：\texttt{output/paper\_figures/fig05\_combined\_visual.\{png,pdf\}}。
\eng{Hybrid+Native-FC} 的 step 序列 panel 可作为 supp 选图补充，不是主文硬门槛。

\eng{Hybrid} 输入变体（V9A）的单独漂移轨迹（原主文 companion，因 proxy 跨输入模式不可横比移至此处）见图~\ref{fig:d-v9a}。

\begin{figure}[ht]
  \centering
  \includegraphics[width=\linewidth]{../../output/paper_figures/fig03s_v9a_trajectory.png}
  \caption{\eng{Hybrid} 输入漂移轨迹（companion；run id: V9A）。artifact / corr vs 训练步数，纵轴尺度仅在 \eng{Hybrid} 自身轨迹内可比，不能与 \(1\times\) 输入变体横向比较。
  \eng{Hybrid} 非单调、晚期压平但平台仍低于早期 checkpoint——输入端证据注入改变漂移形态、不消除先验侵蚀。}
  \label{fig:d-v9a}
\end{figure}

\subsection{零训练融合 baseline}

本节建立一个无需额外训练的简单 baseline：对经典重建（锚）与学习方法输出做事后线性融合，测试学习方法的增益是否超出后处理可得的范围。

\subsubsection{方法}

融合公式为 \(\text{fused}(\lambda) = (1-\lambda) \cdot \text{anchor} + \lambda \cdot \hat{x}_{\text{unet}}\)，\(\lambda \in \{0, 0.1, \ldots, 1.0\}\)。
锚选择 drizzle \(2\times\) mean（248 帧 contour\_refined）或 \eng{TGV}。UNet 预测使用 \eng{Hybrid}（V9A）20K（最保真）和 \eng{Hybrid}（V9A）60K（最锐）。评估指标为 D.3.2 的 fine-window 四指标。

E3 追加了第二个固定验证窗，入口为 \texttt{algos/ep07\_unet\_sr/scripts/v9\_review/run\_fusion\_window2.py}，运行时显式设置 \texttt{CUDA\_VISIBLE\_DEVICES=""} 且只读取缓存全幅预测。
原 D.7 fine-window 为 \(2\times\) 网格 rows 384:518, cols 478:674；第二窗保持同尺寸与同列，向下平移一个窗口高度并留 24\,px 间隔，即 rows 542:676, cols 478:674。
该规则在重选 \(\lambda\) 前固定，用于检查原窗局部选择是否可迁移到另一个非重叠 \eng{ROI}。

\subsubsection{结果}

原 fine-window 上，\eng{TGV} 锚 \(\times\) \eng{Hybrid}-60K（V9A）的低 \(\lambda\) 组合形成一段局部 proxy-frontier。
\(\lambda = 0.2\) 时四指标为 hp\_corr\_input \(=\) 0.963（\eng{TGV} 参照 0.960）、hp\_corr\_tgv \(=\) 0.995、sharp\_p95 \(=\) 0.968（\eng{TGV} 参照 0.959）、lattice \(=\) 0.0108（\eng{TGV} 参照 0.0169）。
这说明在该局部窗口内，少量 \eng{Hybrid}-60K 高频内容可以改善 fine-window proxy 读数；但该读数依赖 \eng{TGV}/highpass 参照与同窗 \(\lambda\) 选择，不是 GT 保真或物理分辨率证据。

第二窗结果位于 \texttt{output/ep07\_v9\_review/fusion\_window2\_metrics.csv} 与 \texttt{fusion\_window2\_summary.md}。
第二窗局部参照为 drizzle \(=\) (hp\_corr\_input 1.000, sharp\_p95 0.427, lattice 0.0036)，\eng{TGV} \(=\) (0.962, 0.491, 0.0160)。
按同一选择规则（\eng{TGV} 锚 \(\times\) \eng{Hybrid}-60K、\(\lambda > 0\)，要求 hp\_corr\_input 与 sharp\_p95 不低于本窗 \eng{TGV}、lattice 不高于本窗 \eng{TGV}，并以 hp\_corr\_input 优先排序），
重选得到 \(\lambda = 0.1\)：hp\_corr\_input \(=\) 0.964、hp\_corr\_tgv \(=\) 0.998、sharp\_p95 \(=\) 0.508、lattice \(=\) 0.0126。
\(\lambda = 0.2\) 也通过同一 gate（0.964 / 0.994 / 0.525 / 0.0100），但因 hp\_corr\_input 略低而排第二。

相比之下，drizzle 锚 \(\times\) \eng{Hybrid} 的任意 \(\lambda\) 都没有形成同类前沿点：drizzle 参照本身保真极高，但锐度始终不足。

\eng{Hybrid+ResObs}（V10）高-\(\lambda\) 工作点的相对位置跨窗保持一致：原窗 \(\lambda = 1.2\) @15K 为 hp\_corr\_input \(=\) 0.922、sharp\_p95 \(=\) 0.987、lattice \(=\) 0.0141；
第二窗为 hp\_corr\_input \(=\) 0.920、sharp\_p95 \(=\) 0.501、lattice \(=\) 0.0130。
两窗中 \eng{Hybrid+ResObs} 均低于对应 \eng{TGV} 的 hp\_corr\_input（0.960/0.962），同时 lattice 低于 \eng{TGV}、sharp\_p95 接近或略高于 \eng{TGV}。
这支持“\eng{Hybrid+ResObs} 是低 grain / 较锐但保真不足的折中点”这一 proxy 位置判断，而不是可认证胜出。

\subsubsection{结论与边界}

第一，零训练线性融合是一个必须报告的强 sanity baseline：在两个非重叠局部窗口中，\eng{TGV} 锚 \(\times\) \eng{Hybrid}-60K 的低 \(\lambda\) 点都能进入相对 \eng{TGV} 的 proxy-frontier 区域。
第二，最优 \(\lambda\) 不完全跨窗稳定（0.2 \(\rightarrow\) 0.1），因此不能把单窗最优点写成稳定算法结论；更稳妥的表述是“低 \(\lambda\) TGV-anchored 融合在局部 proxy 上构成对学习方法的挑战”。

\eng{Hybrid+ResObs} 高-\(\lambda\) sweep 的最佳工作点仍按 \(\lambda = 1.2\) @15K 报告：原窗 hp\_corr\_input \(=\) 0.922、sharp\_p95 \(=\) 0.987、lattice \(=\) 0.0141；第二窗 hp\_corr\_input \(=\) 0.920、sharp\_p95 \(=\) 0.501、lattice \(=\) 0.0130。
它满足“低 grain、锐度不低于 \eng{TGV} 局部读数”的 proxy 判据，但保真仍低于 \eng{TGV}，因此不构成可认证胜出。

边界必须保留：fine-window 与第二窗都是局部 \eng{ROI}，指标依赖 \eng{TGV}/highpass 参照而非 GT；\(\lambda\) 是在同一真实数据窗口上重选，仍有 selection-on-test caveat；
\texttt{sharp\_p95} 不度量轮廓连续性且可受 beading/grain 影响。D.7 的结论只用于 proxy 稳定性与 baseline 压力测试，不能外推为温度计量、物理分辨率或单一方法胜出声明。

\section{复现说明}

\subsection{代码与环境}

\subsubsection{仓库结构}

随附代码包按以下结构组织。

\begin{table}[ht]
  \centering
  \caption{随附代码包目录结构与论文对应}
  \label{tab:repo-structure}
  \small
  \begin{tabular}{p{0.27\linewidth}p{0.46\linewidth}p{0.18\linewidth}}
    \toprule
    目录 & 内容 & 论文对应 \\
    \midrule
    \texttt{core/} & 共享层：数据审计、位移模型、理论界计算、锚点 gate、绘图规范 & \S3 全部标定数字 \\
    \texttt{configs/} & 物理参数单一来源：stage 标定、PSF 标定、噪声底、对齐结果 & \S3.1 标定链 \\
    \texttt{tcforge/} & 物理匹配合成平台（独立包） & \S4.2 \\
    \texttt{algos/ep06\_sr\_poc/} & forward 算子 \(+\) 指标库 & \S3.1 / 附录 A.3 \\
    \texttt{algos/ep07\_unet\_sr/} & 学习方法训练/推理/评估代码 \(+\) 诊断管线 & \S4.3 / \S5 \\
    \texttt{algos/ep10\_*/} & drizzle / MAP-TV / TGV 经典方法 & \S4.1 \\
    \texttt{algos/ep15\_info\_limit/} & FRC、\(\sigma\) 仲裁、去卷积基准 & 附录 A.4 / \S3.1 \\
    \texttt{algos/ep16\_budget\_robustness/} & 帧预算与鲁棒性矩阵 & \S5.4--5.5 \\
    \texttt{scripts/paper\_figures/} & 论文图生产脚本 & F1--F7 \\
    \texttt{paper/reports/}、\texttt{research\_log/} & 正式报告与全程实验日志 & 可追溯性 \\
    \bottomrule
  \end{tabular}
\end{table}

原始数据（\texttt{data/}）和实验产物（\texttt{output/}）不入版本控制，前者需手动放置，后者由脚本重建。

\subsubsection{环境配置}

本项目使用 \eng{UV} 进行依赖管理。根目录 venv（\texttt{uv sync} \(+\) \texttt{uv pip install -e core/}）用于全局脚本、notebook 和数据审计。
每个算法目录（\texttt{algos/xxx/}）是独立项目，拥有自己的 venv 和 lockfile，通过 \texttt{uv sync} 创建。共享代码以 \texttt{pip install -e ../../core} 方式安装。
\eng{TGV} 经典方法因包含 C/CUDA 编译依赖而使用 conda 环境（\texttt{environment.yml}）。各算法环境之间零耦合，可独立运行和删除。

\subsubsection{新机器一键部署}

\begin{verbatim}
git clone <repo> thermal_lift && cd thermal_lift
# place data/data_raw/ (not distributed; see E.3)
uv sync && uv pip install -e core/
uv run python scripts/build_ep01_cache.py       # build frame_audit.csv
uv run python scripts/build_all_notebooks.py --execute  # all artifacts
\end{verbatim}

\subsection{图表重建命令}

下表列出每张论文图表的生成脚本、所需环境和输入依赖。环境标注中 "root" 表示仓库根 venv，"ep07" 等表示对应算法目录内的 venv。

\subsubsection{主文图表}

\begin{table}[ht]
  \centering
  \caption{主文图表生成脚本与环境}
  \label{tab:main-figs}
  \small
  \begin{tabular}{p{0.20\linewidth}p{0.50\linewidth}p{0.09\linewidth}p{0.12\linewidth}}
    \toprule
    图表 & 脚本 & 环境 & 耗时 \\
    \midrule
    F1 系统标定 & \texttt{scripts/paper\_figures/fig01\_system\_calibration.py} & root & 秒级 \\
    F2 FRC 曲线 & \texttt{scripts/paper\_figures/fig02\_frc.py} & root & 秒级 \\
    F3 零空间漂移 & \texttt{algos/ep07\_unet\_sr/scripts/plot\_drift\_trajectories\_paper.py} & ep07 & 秒级 \\
    F4 Pareto 选点 & \texttt{algos/ep07\_unet\_sr/scripts/plot\_checkpoint\_selection.py} & ep07 & 分钟级 \\
    F5 主视觉对比 & \texttt{algos/ep11\_dl\_benchmark/scripts/run\_unified\_harness\_t1\_t2.py} & ep11 & \(\sim\)3\,min \\
    F7 预算鲁棒性 & \texttt{algos/ep16\_budget\_robustness/scripts/run\_ep16\_classical.py --summarize-only} & ep16 & 秒级 / \(\sim\)3.2\,h \\
    \bottomrule
  \end{tabular}
\end{table}

T1/T2 与 F5 同源生成命令：

\begin{verbatim}
cd algos/ep11_dl_benchmark
CUDA_VISIBLE_DEVICES=0 uv run python scripts/run_unified_harness_t1_t2.py \
  --device cuda:0 \
  --workers 4 \
  --output-dir ../../output/ep11_unified_harness
\end{verbatim}

输出位于 \texttt{output/ep11\_unified\_harness/}，包含四个产物文件
\texttt{all\_arm\_metrics.csv}、
\texttt{t1\_metrics.csv}、
\texttt{t2\_metrics.csv} 与
\texttt{run\_manifest.json}；
主视觉图为 \texttt{output/paper\_figures/} 下的 \texttt{fig05\_main\_visual.\{png,pdf\}}。

\subsubsection{V9 Review 诊断管线}

以下脚本在 ep07 环境下依序执行，输出统一到 \texttt{output/ep07\_v9\_review/}：

\begin{verbatim}
uv run python scripts/v9_review/extract_tb_metrics.py             # extract TB scalars
uv run python scripts/v9_review/run_pareto_sweep.py --device cuda # checkpoint inference
uv run python scripts/v9_review/render_comparison_panels.py       # comparison panels
uv run python scripts/v9_review/run_fusion_baseline.py            # fusion sweep
\end{verbatim}

\subsubsection{补充材料图表}

\begin{table}[ht]
  \centering
  \caption{补充材料图表重建脚本}
  \label{tab:supp-figs}
  \small
  \begin{tabular}{p{0.27\linewidth}p{0.50\linewidth}p{0.15\linewidth}}
    \toprule
    资产 & 重建脚本 & 环境 \\
    \midrule
    S-F1 FRC 档案图 & \texttt{scripts/paper\_figures/fig02\_frc.py} & root \\
    S-F2 PSF 证据链 & \texttt{scripts/paper\_figures/figS02\_psf\_evidence.py} & root \\
    S-F9 融合 Pareto & \texttt{scripts/paper\_figures/figS09\_fusion\_pareto.py} & root \\
    S-F10 Hybrid（V9A）演化条带 & \texttt{scripts/paper\_figures/figS10\_v9a\_strip.py} & root \\
    D.2 MAP-TV 基准 & \texttt{algos/ep15\_info\_limit/scripts/run\_m4\_deconv\_anchor.py} & ep15 \\
    \bottomrule
  \end{tabular}
\end{table}

\subsubsection{训练复现}

各变体的完整训练 CLI 存档于 \texttt{algos/ep07\_unet\_sr/scripts/} 下的
\texttt{run\_training.md}、\texttt{run\_v9.md}、\texttt{run\_v10.md} 三个文件。
单个变体 60K 步约需 10--20 小时（取决于 batch size 和 GPU 型号）。

随机性控制：合成数据与 split 全部使用显式 seed（FRC: 42/123/456；EP16 子集: 101/202/303、扰动: 401--403；训练池与 drizzle 变体按 \([\text{seed}, \text{epoch}, \text{scene}]\) 确定性采样）。
长任务输出 run manifest 记录全部配置与结果状态。

\subsection{数据与伦理声明}

\textbf{样品性质。} 论文热像来自实验室采集的工业芯片 \eng{LWIR} 温度矩阵（主 session 248 帧 clean set），全幅 \eng{FOV}（\(640 \times 480\) \eng{LR} / \(1280 \times 960\) \eng{HR}）可直接用于配图，无需 \eng{ROI} 脱敏或型号匿名化。

\textbf{数据分发。} 原始 \eng{TXT} 温度矩阵不随论文公开（实验室数据管理惯例）。读者可在 \eng{TCForge} 合成数据上完整复现训练与评估协议。

\textbf{可公开内容。} \eng{TCForge} 合成平台全部源码、所有标定常数（\(\theta = 47.6°\)、pitch 10\,\(\mu\)m、空间分辨率 20\,\(\mu\)m、PSF \(\sigma\) 区间 [0.2, 0.5]、噪声底 0.0724\,\(^\circ\)C）、评估脚本和 frame\_audit.csv schema 均随代码包提供。

\textbf{诚实声明。} 真实数据结论的第三方复现依赖三层证据：协议（本文）、实验日志（\texttt{research\_log/} 随仓库提供）和中间产物 CSV（\texttt{paper/reports/} 随仓库提供）。


% 补充材料当前不含独立引用；若后续加入 \cite{...}，取消下面两行注释并改用
% `make -C paper zh_conf_supp`（含 bibtex pass）重新编译。
% \bibliographystyle{gbt7714-2015-numeric}
% \bibliography{../bib/refs}

\end{document}
